\phantomsection
\addcontentsline{toc}{subsection}{Inciso c)}
\subsection*{Inciso c)}

% Seccion con nombres de clase largos, que no se parten con guion.
\sloppy

\textit{Analiza: ¿Qué pasa si se quiere añadir una nueva actividad sin modificar el Facade?
¿Cómo lo resuelves con Command o Strategy?}

\vspace{0.4cm}

\subsubsection*{La pregunta tiene una trampa, y conviene desactivarla primero}

Tal como está escrita, la pregunta admite una respuesta tramposa: ``no se puede, porque abrir
la fachada es modificarla''. Es cierto y es inútil. La pregunta con contenido es otra:

\begin{quote}
¿Cuánto cuesta la \emph{primera} actividad nueva, y cuánto cuestan las \emph{siguientes}?
\end{quote}

Esa sí tiene una respuesta medible, y es la que se contesta aquí: la primera cuesta una
modificación, la que abre la fachada; a partir de ahí, cero.

\subsubsection*{Qué pasa si no se hace nada}

Sin ningún arreglo, cada actividad nueva es un método más en \texttt{HomeTheaterFacade}:

\begin{itemize}[leftmargin=1.2cm]
    \item \textbf{Se viola el principio Open/Closed} en el único punto del diseño que debería
          ser estable. La fachada es lo que todos los clientes importan.
    \item \textbf{Cada actividad recompila a todos los clientes}, porque cambia la clase de
          la que dependen.
    \item \textbf{La fachada crece sin límite} y se acerca al objeto dios que el inciso (a)
          ya anunciaba como su forma de degenerar.
    \item \textbf{Nadie fuera del equipo puede añadir actividades.} Si el cine se vende como
          producto, un instalador no puede definir ``modo fiesta'' sin tocar el código
          fuente.
\end{itemize}

Nótese que las tres actividades del enunciado \emph{no} tienen este problema: son el núcleo
de la fachada y está bien que vivan en ella con nombre propio. El problema aparece con la
cuarta, la quinta y la que pida el cliente el mes que viene.

\subsubsection*{Command: la actividad como objeto}

Command convierte cada actividad en un objeto con su propio ciclo de vida. La fachada deja de
conocer las actividades una por una y pasa a conocer el \emph{concepto} de actividad:

{\small
\begin{verbatim}
public interface IActividad
{
    string Nombre { get; }
    void Ejecutar();
    void Deshacer();
}
\end{verbatim}
}

Y la fachada gana, de una sola vez, el registro y el despacho:

{\small
\begin{verbatim}
private readonly Dictionary<string, IActividad> _actividades =
    new(StringComparer.OrdinalIgnoreCase);

private IActividad? _ultima;

public void Registrar(IActividad actividad)
    => _actividades[actividad.Nombre] = actividad;

public void Ejecutar(string nombre)
{
    if (!_actividades.TryGetValue(nombre, out IActividad? actividad))
    {
        throw new KeyNotFoundException(
            "No hay ninguna actividad registrada con el nombre \""
            + nombre + "\"");
    }

    _ultima = actividad;
    actividad.Ejecutar();
}

public void DeshacerUltima()
{
    _ultima?.Deshacer();
    _ultima = null;
}
\end{verbatim}
}

Tres decisiones de esa implementación que no son obvias:

\begin{itemize}[leftmargin=1.2cm]
    \item \textbf{La actividad recibe los componentes, no la fachada.}
          \texttt{ActividadKaraoke} se construye con el amplificador, el dvd, el proyector,
          la pantalla y las luces. Si recibiera la fachada se crearía un ciclo entre las dos
          y, peor, la actividad quedaría limitada a lo que la fachada ya sabe hacer.
    \item \textbf{Un nombre no registrado falla fuerte.} Lanzar es mejor que no hacer nada:
          un nombre mal escrito es un error del cliente, y una sala que se queda a oscuras
          sin decir por qué es peor que una excepción.
    \item \textbf{\texttt{Deshacer} no es decorativo.} Es lo que Command aporta sobre una
          simple lambda, y aquí tiene sentido literal: apagar el cine \emph{es} deshacer la
          actividad que se encendió.
\end{itemize}

\subsubsection*{La medición}

La afirmación ``a partir de ahí, cero'' no se afirma: se mide sobre este repositorio. La
segunda actividad se añadió en un \textit{commit} aislado a propósito:

{\small
\begin{verbatim}
$ git show --stat --format="%h %s" cb5c20d
cb5c20d docs(cp4): agrega una segunda actividad como experimento del inciso (c)

 .../CP4/codigo/ejercicio7/ActividadModoFiesta.cs | 45 ++++++++++++++++++++
 1 file changed, 45 insertions(+)
\end{verbatim}
}

\begin{table}[h!]
    \centering
    \small
    \begin{tabular}{p{5.4cm}p{2.6cm}p{2.6cm}p{3.0cm}}
        \toprule
        \textbf{Cambio} & \textbf{Archivos nuevos} & \textbf{Modificados}
        & \textbf{Sin Command} \\
        \midrule
        Abrir la fachada (una sola vez) & 2 & 1 & --- \\
        \addlinespace
        Primera actividad nueva (karaoke) & 1 & 0 & 1 modificado \\
        \addlinespace
        Segunda actividad nueva (fiesta) & 1 & \textbf{0} & 1 modificado \\
        \bottomrule
    \end{tabular}
\end{table}

El cálculo es sencillo: con $A$ actividades añadidas, el enfoque sin patrón modifica la
fachada $A$ veces y Command una sola. La inversión se amortiza en la primera.

\subsubsection*{Strategy: responde otra pregunta}

Strategy también mete un objeto donde antes había código fijo, de ahí que el enunciado los
nombre juntos. Pero no sirve para lo que la pregunta pide:

{\small
\begin{verbatim}
public interface IPerfilSala
{
    string Nombre { get; }
    int LucesPelicula { get; }
    int VolumenPelicula { get; }
    bool EnvolventePelicula { get; }
}

// en la fachada:
public IPerfilSala Perfil { get; set; } = new PerfilEstandar();

public void VerPelicula(string titulo)
{
    _luces.Atenuar(Perfil.LucesPelicula);
    ...
    _amplificador.PonerVolumen(Perfil.VolumenPelicula);
    ...
}
\end{verbatim}
}

Con \texttt{PerfilNocturno}, la misma película se ve con las luces al 5\%, el volumen a 3 y
sin envolvente. Lo que cambió no es \emph{qué} sabe hacer la sala, sino \emph{cómo} lo hace:
la secuencia sigue teniendo los mismos once pasos.

\begin{table}[h!]
    \centering
    \small
    \begin{tabular}{p{3.4cm}p{5.4cm}p{5.4cm}}
        \toprule
        & \textbf{Command} & \textbf{Strategy} \\
        \midrule
        Qué varía & \emph{qué} se hace: la actividad completa & \emph{cómo} se hace una
                    actividad que ya existe \\
        \addlinespace
        Unidad & una actividad con sus pasos & un conjunto de parámetros o un algoritmo \\
        \addlinespace
        Cuántas conviven & todas a la vez, en un registro & una a la vez, la seleccionada \\
        \addlinespace
        Deshacer & sí, es parte del patrón & no tiene sentido \\
        \addlinespace
        ¿Responde a la pregunta del inciso? & \textbf{Sí} & \textbf{No}, resuelve un
                    problema vecino \\
        \bottomrule
    \end{tabular}
\end{table}

La conclusión no es que Strategy sobre, sino que las dos preguntas son distintas y en este
diseño conviven sin estorbarse: Command abrió el catálogo de actividades y Strategy
parametrizó una de ellas.

\subsubsection*{El precio de abrir la fachada, dicho sin adornos}

Abrir la fachada no sale gratis, y conviene que quede escrito:

\begin{itemize}[leftmargin=1.2cm]
    \item \textbf{Los nombres pasan a ser cadenas.} \texttt{cine.VerPelicula("Blade Runner")}
          lo comprueba el compilador; \texttt{cine.Ejecutar("karaoke")} no. Se cambia
          seguridad en compilación por extensibilidad en ejecución. Es el mismo compromiso
          que el ejercicio 12 discute con los miembros por omisión, y tiene el mismo olor.
    \item \textbf{La fachada se parece más a un despachador.} Llevado al extremo, una
          fachada cuya única operación es \texttt{Ejecutar(string)} ha dejado de simplificar:
          el cliente vuelve a necesitar saber qué nombres existen.
    \item \textbf{Por eso la implementación conserva las dos vías.} Las tres actividades del
          enunciado siguen siendo métodos con nombre y firma propia, que es lo que una
          fachada debe ofrecer; el registro es la puerta para lo que venga después. Elegir
          una sola de las dos vías habría sido más puro y peor.
\end{itemize}

\subsubsection*{Evidencia de ejecución}

Salida literal de \texttt{DemoExtension}:

{\footnotesize
\begin{verbatim}
1. Command: una actividad nueva, escrita sin tocar la fachada
-------------------------------------------------------------
  cine.Registrar(new ActividadKaraoke(...));
  actividades registradas: karaoke

  cine.Ejecutar("karaoke");
    luces         atenuadas al 70%
    pantalla      bajada
    proyector     encendido
    proyector     entrada -> dvd
    amplificador  encendido
    amplificador  entrada -> dvd
    amplificador  modo estereo
    amplificador  volumen 8
    dvd           encendido
    dvd           reproduciendo "Lista de karaoke"

2. Deshacer: lo que Command aporta sobre una simple lambda
----------------------------------------------------------
  cine.DeshacerUltima();
    dvd           detenido
    dvd           apagado
    amplificador  apagado
    proyector     apagado
    pantalla      subida
    luces         al 100%

3. Un nombre que no existe falla fuerte, no en silencio
-------------------------------------------------------
  No hay ninguna actividad registrada con el nombre "ver-la-tele"

4. Strategy: la MISMA actividad, otros parametros
-------------------------------------------------
  perfil estandar (el del inciso b):
    luces         atenuadas al 10%
    pantalla      bajada
    proyector     encendido
    proyector     entrada -> dvd
    proyector     formato 16:9
    amplificador  encendido
    amplificador  entrada -> dvd
    amplificador  modo envolvente
    amplificador  volumen 7
    dvd           encendido
    dvd           reproduciendo "Blade Runner"
    dvd           detenido
    dvd           apagado
    amplificador  apagado
    proyector     apagado
    pantalla      subida
    luces         al 100%

  cine.Perfil = new PerfilNocturno();
    luces         atenuadas al 5%
    pantalla      bajada
    proyector     encendido
    proyector     entrada -> dvd
    proyector     formato 16:9
    amplificador  encendido
    amplificador  entrada -> dvd
    amplificador  modo estereo
    amplificador  volumen 3
    dvd           encendido
    dvd           reproduciendo "Blade Runner"
    dvd           detenido
    dvd           apagado
    amplificador  apagado
    proyector     apagado
    pantalla      subida
    luces         al 100%

5. La segunda actividad nueva ya no costo ninguna modificacion
--------------------------------------------------------------
  actividades registradas: karaoke, modo-fiesta

  cine.Ejecutar("modo-fiesta");
    luces         atenuadas al 80%
    amplificador  encendido
    amplificador  entrada -> dvd
    amplificador  modo envolvente
    amplificador  volumen 10
    dvd           encendido
    dvd           reproduciendo "Lista de fiesta"
    dvd           detenido
    dvd           apagado
    amplificador  apagado
    luces         al 100%

6. Las dos preguntas no son la misma
------------------------------------
  Command cambio QUE sabe hacer la sala: karaoke no existia.
  Strategy cambio COMO hace algo que ya sabia: la pelicula
  sigue teniendo los mismos once pasos, con otros valores.
\end{verbatim}
}

\subsubsection*{Cómo ejecutarlo}

{\small
\begin{verbatim}
cd docs/CP4/codigo/ejercicio7
dotnet run -- extension
\end{verbatim}
}

La salida de \texttt{dotnet run -- demo}, la del inciso (b), no cambió ni un carácter al
introducir Strategy: \texttt{PerfilEstandar} reproduce exactamente los valores que estaban
escritos a mano.

\fussy
